\setvariables[meta]
  [title={Konzept},
   subtitle={Versuchsprotokoll},
   date={2024-10-18},
   author={Marie Kurzhals, Lucas Weckler, Niklas von Hirschfeld},
  ]
\setupinteraction[title={Ausgewogene Ernährung verbunden mit Fitness}]

\setupinteraction[state=start]

\setupheadertexts[\getvariable{meta}{title}][chapter]

\setuphead[chapter][
	header=normal,
	page=no
	before={\testpage[12]\blank[big]},
	after={\blank[2*big]},
]

\starttext

\startfrontmatter
\component c_titlepage
\component c_toc
\stopfrontmatter

\setuppagenumber[number=1]

\startbodymatter

\chapter{Mitagessen - Döner}

\useURL
  [RezeptQuelle]
  [https://www.instagram.com/reel/DAGd3NKI_ep/?igsh=MWsyNjlpc2lrOTVoMw==]
  []
  [source:rezept:doener]

Döner auf Tofu basis. Das Rezept ist von @schmaleschuler\sidenote{\url[RezeptQuelle]} (Auf Insta.)

Zutaten (pro Person):
\startitemize[packed]
\item Tofu (für drei Portionen):
	\startitemize[packed]
	\item 350g geräucherter Todu
	\item 1 Zehe Knoblauch
	\item 1 rote Zwiebel
	\item 1 TL Olivenöl
	\item 2 ELl Soja Sauce
	\item Paprikapulver, Oregano
	\stopitemize
\item Fladenbrot
\item 30ml Joghurt
\item 30g Tomaten
\item 30 Gurke
\item 15g Hirtenkäse
\item Petersilie
\item Salat
\stopitemize

Nährwerte für den Döner komplett ca.:

\startitemize[packed]
\item 647 kcal
\item 33g Protein
\item 80g KH
\item 13g Fett
\stopitemize

\stopbodymatter

\startbackmatter
% \component c_definitions
% \component c_bibliography
\stopbackmatter

\stoptext
